\documentclass[a4paper,10pt]{article}
\usepackage{xeCJK}
\setCJKmainfont{Hiragino Mincho ProN}
\setCJKsansfont{Hiragino Sans}
\setCJKmonofont{Hiragino Sans}

\usepackage[top=12.5mm,bottom=12.5mm,left=13mm,right=13mm]{geometry}
\usepackage{amsmath,amssymb,bm}
\usepackage{multicol}
\usepackage{titlesec}
\usepackage{xcolor}
\usepackage[hidelinks]{hyperref}

\pagestyle{empty}
\setlength{\parindent}{1em}
\setlength{\parskip}{0pt}
\setlength{\columnsep}{5.8mm}
\renewcommand{\baselinestretch}{0.95}

\titleformat{\section}
  {\normalfont\bfseries\normalsize}
  {\thesection.}{0.45em}{}
\titlespacing*{\section}{0pt}{0.45\baselineskip}{0.12\baselineskip}

\newcommand{\E}{\mathbb E}
\newcommand{\Prb}{\mathbb P}
\newcommand{\indep}{\mathrel{\perp\!\!\!\perp}}
\newcommand{\Y}{\bm Y}
\newcommand{\Z}{\bm Z}

\newcommand{\papertitle}{欠測IVと潜在変数モデリングによる非無作為欠測の識別}
\newcommand{\subtitle}{対話用ワーキングドラフト}
\newcommand{\authorblock}{%
  \begin{tabular}{c}
    本多 将大\(^{1,2}\)、星野 崇宏\(^{1,2}\)、大津 泰介\(^{1,3}\)\\[-0.1mm]
    {\scriptsize 1. 慶應義塾大学、2. 理化学研究所、3. London School of Economics and Political Science}
  \end{tabular}
}

\newcommand{\draftprompt}[1]{%
  \noindent{\color{gray}\footnotesize #1}\par\vspace{1.5mm}
}

\begin{document}

\vspace*{-9mm}
\begin{center}
  {\large\bfseries \papertitle\par}
  \vspace{0.2mm}
  {\small \subtitle\par}
  \vspace{0.5mm}
  {\footnotesize \authorblock}
\end{center}
\vspace{-3.5mm}

\footnotesize
\setlength{\baselineskip}{9.7pt}
\setlength{\abovedisplayskip}{1.2pt}
\setlength{\belowdisplayskip}{1.2pt}
\setlength{\abovedisplayshortskip}{0.8pt}
\setlength{\belowdisplayshortskip}{0.8pt}
\setlength{\jot}{0pt}
\begin{multicols}{2}

\section{問題設定}
\draftprompt{最初の一行：観測単位は何か。}
\draftprompt{各単位について、どの変数が存在するか。}
\draftprompt{そのうち何が観測され、何が欠測するか。}
\draftprompt{欠測はどのような仕組みで生じるか。}

% 問題設定は、ユーザーが確定した一文だけを順に追加する。

\section{識別したい対象}
\draftprompt{最終的に知りたい分布、母数、または効果を一文で置く。}

% 問題設定が固まるまでは空欄でもよい。

\columnbreak

\section{識別のアイデア}
\draftprompt{必要な仮定と識別の論理を、確定した順に追加する。}

% 未確定の主張は定理として書かず、仮説またはメモとして扱う。

\section{未解決点}
\draftprompt{まだ決めていない点、反例、既存研究との差を残す。}

% 対話中の保留事項をここに追加する。

\end{multicols}
\end{document}
